\documentclass[10pt,a4paper]{amsart}
\usepackage{amsmath,amssymb,amsthm,graphicx}
\usepackage{verbatim}
\newenvironment{reference}{\comment}{\endcomment}
%\newenvironment{reference}{}{}
\newenvironment{slogan}{\comment}{\endcomment}
\newenvironment{history}{\comment}{\endcomment}

% For commutative diagrams we use Xy-pic
\usepackage[all]{xy}

% We use 2cell for 2-commutative diagrams.
\xyoption{2cell}
% Theorem environments.
%
\theoremstyle{plain}
\newtheorem{theorem}[subsection]{Theorem}
\newtheorem{proposition}[subsection]{Proposition}
\newtheorem{lemma}[subsection]{Lemma}

\theoremstyle{definition}
\newtheorem{definition}[subsection]{Definition}
\newtheorem{example}[subsection]{Example}
\newtheorem{exercise}[subsection]{Exercise}
\newtheorem{situation}[subsection]{Situation}

\theoremstyle{remark}
\newtheorem{remark}[subsection]{Remark}
\newtheorem{remarks}[subsection]{Remarks}

\numberwithin{equation}{subsection}

% Macros
%
\def\lim{\mathop{\mathrm{lim}}\nolimits}
\def\colim{\mathop{\mathrm{colim}}\nolimits}
\def\Spec{\mathop{\mathrm{Spec}}}
\def\Hom{\mathop{\mathrm{Hom}}\nolimits}
\def\Ext{\mathop{\mathrm{Ext}}\nolimits}
\def\SheafHom{\mathop{\mathcal{H}\!\mathit{om}}\nolimits}
\def\SheafExt{\mathop{\mathcal{E}\!\mathit{xt}}\nolimits}
\def\Sch{\mathit{Sch}}
\def\Mor{\mathop{\mathrm{Mor}}\nolimits}
\def\Ob{\mathop{\mathrm{Ob}}\nolimits}
\def\Sh{\mathop{\mathit{Sh}}\nolimits}
\def\NL{\mathop{N\!L}\nolimits}
\def\CH{\mathop{\mathrm{CH}}\nolimits}
\def\proetale{{pro\text{-}\acute{e}tale}}
\def\etale{{\acute{e}tale}}
\def\QCoh{\mathit{QCoh}}
\def\Ker{\mathop{\mathrm{Ker}}}
\def\Im{\mathop{\mathrm{Im}}}
\def\Coker{\mathop{\mathrm{Coker}}}
\def\Coim{\mathop{\mathrm{Coim}}}

% Boxtimes
%
\DeclareMathSymbol{\boxtimes}{\mathbin}{AMSa}{"02}

%
% Macros for moduli stacks/spaces
%
\def\QCohstack{\mathcal{QC}\!\mathit{oh}}
\def\Cohstack{\mathcal{C}\!\mathit{oh}}
\def\Spacesstack{\mathcal{S}\!\mathit{paces}}
\def\Quotfunctor{\mathrm{Quot}}
\def\Hilbfunctor{\mathrm{Hilb}}
\def\Curvesstack{\mathcal{C}\!\mathit{urves}}
\def\Polarizedstack{\mathcal{P}\!\mathit{olarized}}
\def\Complexesstack{\mathcal{C}\!\mathit{omplexes}}
% \Pic is the operator that assigns to X its picard group, usage \Pic(X)
% \Picardstack_{X/B} denotes the Picard stack of X over B
% \Picardfunctor_{X/B} denotes the Picard functor of X over B
\def\Pic{\mathop{\mathrm{Pic}}\nolimits}
\def\Picardstack{\mathcal{P}\!\mathit{ic}}
\def\Picardfunctor{\mathrm{Pic}}
\def\Deformationcategory{\mathcal{D}\!\mathit{ef}}
\usepackage{cleveref}
\crefname{originallemma}{original lemma}{original lemmas}
\Crefname{originallemma}{Original lemma}{Original lemmas}
\makeatletter
\newlabel{chow-lemma-gysin-fundamental}{{0FBM}{source}{}{sourcepointer-1}{}}
\newlabel{chow-lemma-gysin-fundamental@cref}{{[originallemma][0][]0FBM}{[1][0][]source}{}{}{}}
\makeatother
\crefname{originallemma}{original lemma}{original lemmas}
\Crefname{originallemma}{Original lemma}{Original lemmas}
\makeatletter
\newlabel{chow-lemma-intersect-regularly-embedded}{{0FFE}{source}{}{sourcepointer-2}{}}
\newlabel{chow-lemma-intersect-regularly-embedded@cref}{{[originallemma][0][]0FFE}{[1][0][]source}{}{}{}}
\makeatother
\crefname{originalreference}{original reference}{original references}
\Crefname{originalreference}{Original reference}{Original references}
\makeatletter
\newlabel{definition-weil-cohomology-theory-classical}{{0FGV}{source}{}{sourcepointer-3}{}}
\newlabel{definition-weil-cohomology-theory-classical@cref}{{[originalreference][0][]0FGV}{[1][0][]source}{}{}{}}
\makeatother
\crefname{originalreference}{original reference}{original references}
\Crefname{originalreference}{Original reference}{Original references}
\makeatletter
\newlabel{lemma-class-diagonal-classical}{{0FGZ}{source}{}{sourcepointer-4}{}}
\newlabel{lemma-class-diagonal-classical@cref}{{[originalreference][0][]0FGZ}{[1][0][]source}{}{}{}}
\makeatother
\crefname{originalreference}{original reference}{original references}
\Crefname{originalreference}{Original reference}{Original references}
\makeatletter
\newlabel{lemma-degrees-cycles-classical}{{0FGW}{source}{}{sourcepointer-5}{}}
\newlabel{lemma-degrees-cycles-classical@cref}{{[originalreference][0][]0FGW}{[1][0][]source}{}{}{}}
\makeatother
\crefname{originallemma}{original lemma}{original lemmas}
\Crefname{originallemma}{Original lemma}{Original lemmas}
\makeatletter
\newlabel{morphisms-lemma-differentials-diagonal}{{08S2}{source}{}{sourcepointer-6}{}}
\newlabel{morphisms-lemma-differentials-diagonal@cref}{{[originallemma][0][]08S2}{[1][0][]source}{}{}{}}
\makeatother
\title{The self intersection of the diagonal for a classical Weil cohomology theory}
\author{The Stacks Project Authors}
\date{stacks-project:weil.tex:local-snapshot-20261012; GNU FDL 1.2 or later}
\begin{document}
\maketitle
\noindent\textit{Editorial note.} All delivered mathematical text is the author's own, copied from the pinned Stacks Project chapter file weil.tex; no statement or proof step was rewritten, reordered or removed.\par\medskip
\noindent\textit{Editorial note.} The counted proof is the author's complete argument: the identification of the self intersection with the degree of the top Chern class of the tangent sheaf, and the computation through the degree map and the Kunneth decomposition of the class of the diagonal.\par\medskip
\noindent\textit{Editorial note.} The label lemma-square-diagonal-classical names the counted statement, and the proof environment immediately following it in the source is delivered verbatim as its proof.\par\medskip
\noindent\textit{Editorial note.} The counted statement is carried by the author's own lemma environment in the chapter source.\par\medskip
\noindent\textit{Editorial note.} Pre-open gate: python operations/agent-authoring/tools/precard\_gate.py --source weil.tex --label lemma-square-diagonal-classical --start 1801 --end 1849 --local-id 28139 returned RESULT=OK (L28139 free, occupancy FREE, public interval clear). Pre-assembly tag self-check: canonical tag resolved and not present in reference/public-items.jsonl.\par\medskip
\noindent\textit{Editorial note.} The statement and proof are copied from the pinned chapter file; the cited inputs are kept as external pointers with their canonical Stacks tags.\par\medskip
\noindent\textit{Allowed standard prerequisites (editorial).}\begin{itemize}
\item The self intersection of the diagonal through the Gysin map and the top Chern class. Chow Homology, the lemma expressing the self intersection of the diagonal through the Gysin map, and the lemma identifying the Gysin map with capping by the top Chern class of the normal sheaf. Produces the identification of the self intersection with the degree of the top Chern class of the tangent sheaf.
\item The conormal sheaf of the diagonal is the sheaf of differentials. Morphisms, the lemma identifying the conormal sheaf of the diagonal with the sheaf of differentials. Identifies the normal sheaf of the diagonal with the tangent sheaf of the variety.
\item The degree map and the Kunneth decomposition of the class of the diagonal. Same chapter, the classical version of the description of degrees of cycles through the degree map, the lemmas on the pushforward of the unit and on pushforwards, and the classical version of the description of the class of the diagonal. Computes the same intersection product through the cohomological degree map and the Kunneth components of the diagonal class.
\end{itemize}\medskip
\ifdefined\hypertarget\hypertarget{sourcepointer-1}{}\fi
\noindent\textit{Source reference.} Original lemma 0FBM: sources/stacks-chow/tags.txt, source lines 19152--19152. Prerequisite pointer invoked at proof line(s) 1820; it names a standard fact of a sibling chapter that the counted argument uses, and every step of the counted proof built on it is retained verbatim in the delivered body.\par\medskip
\ifdefined\hypertarget\hypertarget{sourcepointer-2}{}\fi
\noindent\textit{Source reference.} Original lemma 0FFE: sources/stacks-chow/tags.txt, source lines 19284--19284. Prerequisite pointer invoked at proof line(s) 1815; it names a standard fact of a sibling chapter that the counted argument uses, and every step of the counted proof built on it is retained verbatim in the delivered body.\par\medskip
\ifdefined\hypertarget\hypertarget{sourcepointer-3}{}\fi
\noindent\textit{Source reference.} Original reference 0FGV: sources/stacks-weil/pinned-original-source.tex, source lines 1682--1682. Prerequisite pointer invoked at proof line(s) 1804; it names a standard fact of the same chapter that the counted argument uses, and every step of the counted proof built on it is retained verbatim in the delivered body.\par\medskip
\ifdefined\hypertarget\hypertarget{sourcepointer-4}{}\fi
\noindent\textit{Source reference.} Original reference 0FGZ: sources/stacks-weil/pinned-original-source.tex, source lines 1753--1753. Prerequisite pointer invoked at proof line(s) 1841; it names a standard fact of the same chapter that the counted argument uses, and every step of the counted proof built on it is retained verbatim in the delivered body.\par\medskip
\ifdefined\hypertarget\hypertarget{sourcepointer-5}{}\fi
\noindent\textit{Source reference.} Original reference 0FGW: sources/stacks-weil/pinned-original-source.tex, source lines 1695--1695. Prerequisite pointer invoked at proof line(s) 1828; it names a standard fact of the same chapter that the counted argument uses, and every step of the counted proof built on it is retained verbatim in the delivered body.\par\medskip
\ifdefined\hypertarget\hypertarget{sourcepointer-6}{}\fi
\noindent\textit{Source reference.} Original lemma 08S2: sources/stacks-chow/tags.txt, source lines 11115--11115. Prerequisite pointer invoked at proof line(s) 1822; it names a standard fact of a sibling chapter that the counted argument uses, and every step of the counted proof built on it is retained verbatim in the delivered body.\par\medskip
\begin{lemma}
\label{lemma-square-diagonal-classical}
Let $H^*$ be a classical Weil cohomology theory
(Definition \ref{definition-weil-cohomology-theory-classical}).
Let $X$ be a smooth projective variety. We have
$$
\sum\nolimits_{i = 0, \ldots, 2\dim(X)} (-1)^i\dim_F H^i(X) =
\deg([\Delta] \cdot [\Delta]) = \deg(c_d(\mathcal{T}_X) \cap [X])
$$
\end{lemma}

\begin{proof}
Equality on the right. We have
$[\Delta] \cdot [\Delta] = \Delta_*(\Delta^![\Delta])$
(Chow Homology, Lemma \ref{chow-lemma-intersect-regularly-embedded}).
Since $\Delta_*$ preserves degrees of $0$-cycles it suffices to compute
the degree of $\Delta^![\Delta]$. The class $\Delta^![\Delta]$ is given
by capping $[\Delta]$ with
the top Chern class of the normal sheaf of $\Delta \subset X \times X$
(Chow Homology, Lemma \ref{chow-lemma-gysin-fundamental}).
Since the conormal sheaf of $\Delta$ is $\Omega_{X/k}$
(Morphisms, Lemma \ref{morphisms-lemma-differentials-diagonal})
we see that the normal sheaf is equal to the tangent sheaf
$\mathcal{T}_X = \SheafHom_{\mathcal{O}_X}(\Omega_{X/k}, \mathcal{O}_X)$
as desired.

\medskip\noindent
Equality on the left. By Lemma \ref{lemma-degrees-cycles-classical} we have
\begin{align*}
\deg([\Delta] \cdot [\Delta])
& =
\int_{X \times X} \gamma([\Delta]) \cup \gamma([\Delta]) \\
& =
\int_{X \times X} \Delta_*1 \cup \gamma([\Delta]) \\
& =
\int_{X \times X} \Delta_*(\Delta^*\gamma([\Delta])) \\
& =
\int_X \Delta^*\gamma([\Delta])
\end{align*}
Write $\gamma([\Delta]) = \sum  e_{i, j} \otimes e'_{2d - i , j}$
as in Lemma \ref{lemma-class-diagonal-classical}.
Recalling that $\Delta^*$ is given by cup product we obtain
$$
\int_X \sum\nolimits_{i, j} e_{i, j} \cup e'_{2d - i, j} =
\sum\nolimits_{i, j} \int_X e_{i, j} \cup e'_{2d - i, j} =
\sum\nolimits_{i, j} (-1)^i = \sum (-1)^i\beta_i
$$
as desired.
\end{proof}

\end{document}
